% Generated draft: requires mathematical and source review.
\documentclass[11pt]{article}
\usepackage[margin=25mm]{geometry}
\usepackage{fontspec}
\setmainfont{Latin Modern Roman}
\usepackage{amsmath,amssymb,mathtools,unicode-math}
\setmathfont{Latin Modern Math}
\usepackage{enumitem,needspace,xurl,hyperref}
\setlength{\parindent}{0pt}
\setlength{\parskip}{5pt}
\setlength{\emergencystretch}{3em}
\newcommand{\N}{\mathbb{N}}
\newcommand{\Z}{\mathbb{Z}}
\newcommand{\Q}{\mathbb{Q}}
\newcommand{\R}{\mathbb{R}}
\newcommand{\C}{\mathbb{C}}
\newcommand{\F}{\mathbb{F}}
\newcommand{\A}{\mathbb{A}}
\newcommand{\PP}{\mathbb{P}}
\newcommand{\E}{\mathbb{E}}
\newcommand{\eps}{\varepsilon}
\newcommand{\dd}{\,\mathrm{d}}
\begin{document}
\section*{UnsolvedMath ID: 700493}
\textbf{Source problem code:} catalogue entry
\subsection{Definitions and mathematical statement}
Let \(G\) be a finite group and let \(T\in\operatorname{Syl}_2(G)\), so that \(T\) has order equal to the largest power of \(2\) dividing \(|G|\). For \(B\leq G\), define its normalizer by
\[
N_G(B)=\{g\in G:gBg^{-1}=B\}.
\]
Write \(B\operatorname{char}T\) when every automorphism of \(T\) preserves \(B\). Set
\[
H_T=\left\langle N_G(B)\;\middle|\;1\neq B\leq T,\ B\operatorname{char}T\right\rangle.
\]
Here the brackets denote the subgroup generated by all the indicated normalizers, and the subgroup generated by an empty family is \(1\). The subgroup \(B=T\) is included when \(T\neq1\); the source does not require \(B\) to be proper in \(T\).

Use the standard convention that a \(2\)-local subgroup of \(G\) is a subgroup \(N_G(Q)\) for some nonidentity \(2\)-subgroup \(Q\leq G\); it is not additionally required to be proper in \(G\). For a finite group \(M\), let \(O_2(M)\) denote its largest normal \(2\)-subgroup. Let \(F(M)\) be its largest nilpotent normal subgroup and let \(E(M)\) be the subgroup generated by its components, where a component is a subnormal quasisimple subgroup. A group \(L\) is quasisimple if \(L=[L,L]\) and \(L/Z(L)\) is nonabelian simple. The generalized Fitting subgroup is
\[
F^*(M)=F(M)E(M).
\]

\textbf{Classification problem.} Describe, up to isomorphism, all finite groups \(G\) admitting \(T\in\operatorname{Syl}_2(G)\) such that
\[
H_T<G
\]
and
\[
\text{for every }M\leq G,\qquad
\left(T\leq M\ \text{and}\ M=N_G(Q)\text{ for some }1\neq Q\leq G\text{ of }2\text{-power order}\right)
\Longrightarrow F^*(M)=O_2(M).
\]
The source does not explicitly assume \(T\neq1\), and no such assumption is added here. The definition of a \(2\)-local subgroup above and the interpretation of the malformed generated-subgroup expression require confirmation against the original source. The supplied record labels the problem open; that label is supplied catalog information, not an independently verified mathematical assessment.
\subsection{Short English statement}
Classify the finite groups whose Sylow \(2\)-subgroup \(T\) has its nonidentity characteristic-subgroup normalizers generating a proper subgroup, and for which every \(2\)-local subgroup containing \(T\) has generalized Fitting subgroup equal to its largest normal \(2\)-subgroup. The draft retains the source's lack of an explicit assumption that \(T\neq1\).
\subsection{Sources}
\begin{itemize}
\item Local UnsolvedMath registry, permanent ID \(700493\), source code \texttt{catalogue entry}, title \emph{Kourovka Notebook Problem 8.43}; supplied complete source record, upstream revision \texttt{d371ee917ac9647fbb06cd6fc529deec6f3dda82}.
\item Supplied bibliographic attribution: \emph{The Kourovka Notebook (Unsolved Problems in Group Theory)}, problem 8.43, 8th issue (1982), as included in the cited 21st-issue edition (2026), arXiv identifier \texttt{1401.0300v48}. The supplied set description names E. I. Khukhro and V. D. Mazurov as editors. This reference has not been independently consulted.
\item The record's proposer field and background attribute the problem to A. A. Makhn\"ev, with proposed year 1982; the supplied statement begins with the attribution ``F. Timmesfeld.'' Both attributions are retained without resolving their respective roles.
\item The registry supplies the status \texttt{open}. No independent literature verification or literature-check date is supplied.
\end{itemize}
\end{document}
